
\chapter{Introduction: Intercultural Communication in a Time of Ecological Crisis}
\lhead{\textit{Introduction: Intercultural Communication in a Time of Ecological Crisis}}

\textit{\textbf{Chapter Summary:} This chapter introduces ecological issues as a topic for intercultural communication research. A literature review outlines the need for a critical research program addressing the theme. The research problem is then stated and the aims and structure of the dissertation are outlined.}
\newline

In 2017, more than 15,000 scientists from 184 countries issued a ``warning to humanity'' that Earth's ecosystems are being pushed beyond their capacities to support life. Human-induced changes to the global environment are now so significant that some scientists argue the Earth has entered a new geological epoch, the Anthropocene \citep{Lewis2015a, Waters2016a}. With rapid industrialization, the pace of these changes is likely to accelerate in the coming decades. We are amid the most rapid period of natural resource development and infrastructure expansion in human history. By 2030, trillions of dollars need to be invested in basic infrastructure simply to meet UN Development Goals \citep{UNCTAD2014}. The amount of minerals, ores, fossil fuels, and biomass consumed globally is projected to triple by 2050 \citep{NationalIntelligenceCouncil2013a}. The convergence of ecological pressures and rapid resource development raises unprecedented challenges.

No doubt, the challenges are be technological. However, there will perhaps be even greater challenges related to communication and cooperation among diverse groups of people. Confronting climate extremes, resource scarcity, and other environmental changes will require mobilization of all segments of society. Solutions demand a range of perspectives and know-how. Cultural perspectives are needed to understand the current situation as well as to draw on new ideas and ways of living. Communication, both within and across borders, is essential for these collaborative efforts. In short, the ecological crisis demands intercultural communication.

\section{The Role for Intercultural Communication}

In many respects, an intercultural perspective on environmental issues is not new. There have long been calls for unified global action to confront environmental threats. One could even view global environmentalism as a case study in awareness and consensus building across borders. Overcoming national and cultural differences for the sake of the planet has been a key theme of the environmental movement since at least the 1960-70s. As \cite*{Jasanoff2004a} points out, environmentalism\textemdash alongside nuclear non-proliferation, human rights, and anti-terrorism\textemdash is one of few cases of global ``norms-making'' and supranational governance (32). There have even been successful international agreements in the face of ecological threats. The Montreal Protocol and the reduction of ozone depleting substances is often cited as an example of unified international cooperation \citep{EuropeanCommission2007a}.

While, on the one hand, environmentalism is a case study in international cooperation, on the other, humanity is failing to address the most pressing environmental issues. Despite hundreds of international treaties having been signed in the last half-century \citep{Mitchella}, evidence suggests there is continued degradation of ecological processes essential to support life on the planet \citep{Steffen2015c, Steffen2015b}. One could point to climate change as just one crucial area where international cooperation has failed. 

One might question whether humanity's failure to confront environmental issues is a problem of intercultural communication. In the case of climate change, plausible reasons for not reducing emissions are a combination of structural, economic, technological, and political factors. While intercultural misunderstandings may play a role, it seems doubtful to attribute global policy inaction to cultural differences. Realpolitik and apprehension about economic implications, for example, are more likely obstacles preventing emissions reduction agreements among the some 195 nations involved in climate negotiations \citep{Stern2018a}. 

Yet, international climate agreements are just one aspect of the issue. Moving from the global policy arena to everyday life, intercultural dynamics come into clearer focus. Environmental issues are felt by people, their families, and local communities. It is in everyday life where the human consequences are experienced and disagreements take place. Consider the Dakota Access Pipeline. This was a controversial \$3.8 billion project intended to transfer shale oil from the Northern Plains of the Unites States to the industrial heartland. In August 2016, indigenous protesters chained themselves to heavy machinery in North Dakota to block its construction. In the following months, viewers worldwide saw protesters arrested, attack dogs unleashed, encampments bulldozed, and the heavily armed national guard march in to face off against pipeline protestors. There were many \textit{social} factors at play (political, economic, legal, etc.) in this pipeline debate. However, as will be argued in a subsequent chapter, resistance to this project was fundamentally a \textit{cultural} act. In other words, resistance to this pipeline stemmed from shared history, identities, worldviews, and values. Culture is not only central to understanding this pipeline protest, but the countless other cases worldwide where infrastructure and natural resources are flashpoints for misunderstanding and conflict.	

When it comes to the environment and natural resources, there is a disorienting array of perspectives and interests. Culture and identity are fused with political and economic realities. In the everyday communities where people live and work, the environment is not an abstraction, nor is it reducible to a biophysical entity for detached scientific observation. The environment is the source of health and well-being. It is also the subjective \textit{Umwelt}\footnote{At occasional points in the dissertation original German terms are used. This is often the case for terms that have a specific meanings in the phenomenological tradition (e.g. \textit{Lebenswelt}). \textit{Umwelt} emphasizes the environment as subjectively experienced and invokes the semiotic theories of Jakob von Uexk\"ull and Thomas A. Sebeok (discussed in later sections).} consisting of places and relationships that have cultural and spiritual meanings. At the same time, global political and structural factors remain crucial determinants of the fate of these places and relationships. In short, the topic of environmental change is broad and involves many complex questions and factors. It is precisely in understanding and sorting through such complexities that intercultural communication (ICC) research can play a role. 

\section{Research Gap}

Despite the important role it \textit{could} play, there remains a lack of research looking at intercultural dynamics of environmental issues. To be sure, environmental research is taking place in disciplines related to ICC (e.g. anthropology and communications studies). However, this research does not necessarily address intercultural interactions. Conversely, research that is focused on intercultural communication rarely addresses ecological issues.	

The relation between human culture and the environment has been a topic of anthropological studies since the 1960s \citep[154-157]{Kottak1999a,Perry2003a} coinciding with the emergence of ecological anthropology and cultural ecology \citep{Steward1972a}. However, ecologically-themed anthropology research is often cross-cultural and ethnographic rather than explicitly intercultural. The paucity of intercultural themes in ecological anthropology scholarship is evident by simply looking at published research topics. Searching the \textit{Journal of Ecological Anthropology} for the term ``intercultural'' yields only 2 articles; the same search in \textit{Ecological and Environmental Anthropology} returns none.  These results are perhaps unsurprising given anthropology's drift away from ICC \citep{Leeds-Hurwitz1987a} but, nonetheless, underscore how there is a rich body of research that is cultural but not intercultural. 

One could point to communication studies as a more likely source of intercultural-environmental research. The subfield of environmental communication \citep{Flor2004a, Corbett2006a} spans rhetoric and discourse, media and journalism, public participation, advocacy campaigns, risk communication, and representations of nature in popular culture \citep{Cox2010a}. Yet, communication between cultures remains a relatively rare theme in this subfield. An ``intercultural'' search in the journal \textit{Environmental Communication} returns 11 results, only one of which contains ``intercultural'' as a keyword. Similar conclusions can be drawn when expanding the search to encompass a range of fields related to communication and linguistics. In a meta-analysis of the database \textit{Communication and Mass Media Complete}, \cite*{Mendoza2016a} found that nearly 90 percent of articles with ``environment'' and ``ecology'' as keywords used these terms analogically (e.g. social environment) rather than in reference to nature or the biosphere (3). Evidently, the same authors found it to be even more rare for research from within intercultural communication studies to focus on environmental topics.

\subsection{The Case for a Unifying Framework}

Some of the most well-known comparative frameworks for studying cultures identify the human relation to nature as fundamental to human cultures. \cite{Kluckhohn1961} include relationship to the environment as one of six dimensions with which a society can be categorized. In the Schwartz Value Survey, environmental protection is also considered (under ``Universalism'') as a factor upon which to compare national cultures \citep{Sagiv2000}. Nonetheless, with the possible exception of \cite*{Mendoza2016a} (discussed in the next chapter), the human relationship to nature has not been systematically taken up by intercultural researchers. While there are some studies addressing related topics, a more comprehensive framework is lacking. There are numerous high-level policy materials combining the phrases ``intercultural dialogue'' and ``sustainable development,'' but generally these have not been part of critical research programs. A look at where intercultural communication and the environment is, indeed, being researched requires a rather broad review across several themes and disciplines. In many cases, research only touches on the intercultural aspects of environmental issues.

One prominent theme is community and professional education. For example, looking at educational services for sustainable development in Sub-Saharan Africa, \cite*{Evani2016a} draw from intercultural communication research to analyze conflicting paradigms and goals. Intercultural communication has also been considered as part of an interdisciplinary framework for sustainable development education at the national level \citep{Volodymyr2017a}. In the professional context, \cite*{Merfeld2017a} assess whether intercultural competence was developed through study abroad programs for civil/environmental engineering students.

Another area of practical importance is intercultural risk communication in the face of natural disaster management and prevention. Using the example of the 2011 earthquake in Japan, \cite*{Neuliep2017a} points out that responses to natural disasters are shaped by a culture's value orientations. Studies addressing disasters from an ethnographic and cultural standpoint have focused on community and psychological resilience \citep[e.g.][]{Marsella2004a}. Given the increasingly international scale of natural disasters (both in terms of the impact of events and humanitarian responses), further research is needed that focuses on disaster communication across national and cultural boundaries.

Various analyses of cultural aspects of climate change also offer promising approaches for further intercultural research. \cite*{Krøvel2011a} looks at how climate change media reporting depends on culturally variant frames. Such frame analyses are often cross-cultural. For example, \cite{Xie2015a} does a comparison of climate change framing in US and Chinese newspapers. Frame analysis is crucial because, as \cite*{Rudiak-Gould2013a} points out, climate change skepticism stems from cultural and ideological factors rather than universals. A unifying premise in these studies is that environmental discourses are a reflection mental models and cognition \citep{Lakoff2010a}. Further research might expand frame analysis to a wider range of environmental issues or, more ambitiously, examine the interface between culture, cognition, and ecology.

Although all of the work mentioned above is important in its own right, there lacks unifying themes and methods that root ecological issues within intercultural communication studies. One possible exception is the field of stakeholder relations or stakeholder analysis, which has aims and methods that overlap with intercultural communication. However, as discussed below, the instrumental nature of this concept and its proximity to strategic corporate communications could be problematic from the perspective of both critical intercultural communication and ecological conservation. 

\subsection{Beyond Stakeholders}

The term \textit{stakeholder} has become common in natural resource and environmental management. Stakeholders are identified as distinctive interest groups that are affected by projects and policies related to natural resources and conservation \citep{Reed2009}. While there is no cross-cutting definition of what constitutes a stakeholder in a given situation \citep{Billgren2008}, cultural identity is, no doubt, a key factor. Insights from intercultural communication occur when stakeholders consist of people from different cultures, which is practically an inevitability in the modern world. 

Stakeholder analysis has roots in policy and business management; the former being concerned with power and influence in the policy process, the latter with threats and opportunities that could affect the success of the firm \citep{Varvasovszky2000}. In many sectors, notably the natural resource sector, stakeholder analysis has included systematic relationship mapping and charting the interest/influence of different actors \citep[e.g.][]{Lindenberg1981}. Such approaches to the analysis and management of stakeholders, can be described as instrumental, meaning they are intended to influence and achieve desired outcomes. Even studies with an intercultural focus could be described as  instrumental such as, for instance, \cite*{Wang2014a} who use an intercultural competence model to assess public relations management in the Peruvian mining industry. By contrast, normative approaches are also found in natural resource and environmental management literature, often under the banner of stakeholder participation or communication. Normative approaches employ notions of justice, democracy, or morality to assess legitimacy among stakeholders \citep[1935-36]{Reed2009}. Such approaches might stress stakeholder participation, equity, and involvement of marginalized groups in decision making processes \citep{Johnson2004}.  

This dissertation proposes that overcoming environmental challenges and conflicts requires a move beyond the notion of stakeholders, towards more in-depth understandings of communication itself. This is not to suggest there are not merits to stakeholder analysis as a discipline and practice. However, stakeholder approaches can be problematic when it comes to intercultural communication and environmental debates. 

The very definition of a stakeholder as anyone affected by a decision \citep{Freeman1983}, is itself problematic from an intercultural standpoint. The natural world is a source of cultural identity. People with historical, cultural, and spiritual relationships to landscapes and lifeforms are more than stakeholders to be considered alongside institutions, corporate entities, and others whose interests are often material and bureaucratic. A cultural relationship to the natural world is one of dwelling, care, and meaning. Understanding conflicts related to natural resources requires a new paradigm of cultural analysis that goes beyond most mainstream stakeholder methodologies.

While the need to change paradigms is most evident with respect to instrumental approaches to stakeholder analysis, it is also borne out of inadequacies in normative frameworks. Normative stakeholder communication theory is often premised on \citeauthor{Habermas1984a}' (\citeyear{Habermas1984a}) communicative rationality, which aims for rational agreement through dialogue to establish shared understanding and consensus. This aim is underpinned by the premise that transparent and clear language, as opposed to force or manipulation, has the ability to generate consensus. This aim may seem amenable to intercultural understanding but the issue here, as \cite{Czobor-Lupp2008} points out, is the assumption of linguistic  clarity, transparency, and rationality. Language can, of course, be all of those things. However, language\textemdash particularly when imbued with cultural meanings\textemdash is also  aesthetic, rhetorical, and metaphorical \citep[430]{Czobor-Lupp2008}. In short, normative stakeholder approaches fail to address the complexity and depth of human cultures and communication. 

An intercultural perspective reminds us that verbal communication is just one aspect of communication. Gestures, expressions, paralanguage, and nonverbal communication more broadly, are all inseparable from meaning and understanding. Moreover, as will be elaborated in subsequent chapters, thought and communication are largely unconscious. For this reason, \cite{Lakoff2010a} justifiably claims that an Enlightenment ideal of language and reason is a barrier to understanding why people hold certain views about environmental issues. 

A final reason for the need to overcome a stakeholder approach is the status of nature itself within these frameworks. \cite*{Starik1995} reminds us that most definitions of stakeholders consist only of human entities. The idea of nature and other lifeforms as stakeholders is often overlooked. This points to further flaws in normative frameworks in that notions of justice and equity that underpin stakeholder theory do not necessarily translate into an environmental ethic. To put it bluntly, achieving goals of consensus and participation do not guarantee that humans are not destroying ecosystems. 

\subsection{Addressing the Gap}

Considered as a whole, the literature points to shortcomings in terms of both depth and breadth. The first shortcoming (depth) refers to a lack of guiding conceptual principles. The few studies that do touch on ecology and intercultural communication, lack an explicit discussion of the assumptions underlying the research. Likewise, with its basis in communicative rationality and strategic communication, stakeholder-related research often does not hold up in the face of conceptual challenges posed by intercultural communication. 

What's needed is a conceptual framework that can enable paradigm change at the interface of culture, communication, and ecology. Explicit conceptual and philosophical premises would underscore the complexity and richness of the subject matter and unify disperse topics. \cite*{Busch2017a} likewise suggest that bridging intercultural communication and global sustainability will require intercultural communication to be more explicit about the normative concepts underlying the research. This need for normative concepts points to a broader issue; namely, the planetary ecological crisis raises questions that current paradigms in culture and communication studies may not be suited to address. The need for conceptual re-examination is also implied in \citeauthor{Mendoza2016a}'s (\citeyear{Mendoza2016a}) call for an ``ecological turn'' in intercultural communication research. 

The second shortcoming (breadth) refers to the range and scope of analysis, both geographically and thematically. Much research that does touch on interculturality and the environment is in case study format, confined to single communities, events, or national cultures. Of course, there is great value in these studies. Local field research is particularly crucial and, by its nature, will generally be geographically focused to specific regions and communities. However, intercultural research has an imperative to consider communicative interactions across multiple geographic scales. Principles that apply across borders and speech communities can establish a reference point for further, perhaps more localized work. A related aspect of breadth concerns thematic levels of analysis. Both intercultural and environmental topics are complex and interdisciplinary. Taken together, the various studies highlight the many factors at play: political, economic, cultural, cognitive, etc. Few studies, however, integrate these multiple factors. 

It is possible to address the issues of depth and breath simultaneously. Theories and concepts guiding research would need to integrate multiple levels of analysis across geographic scales. In short, the approach would be \textit{multilevel}, interdisciplinary, and holistic. The requirement for holistic, multilevel approaches is not new to intercultural communication research. For example, the social ecological model \citep{Brofenbrenner1977a,Brofenbrenner1979a} and meso analysis \citep{Rousseau1944a} have been adopted to study intercultural interactions. Such approaches address the theme of geographic scope by integrating the individual, household, community, institution, state, and global levels. However, these approaches do not necessarily integrate multiple factors in an interdisciplinary manner such as, for example, the political, economic, and linguistic aspects. Moreover, the application of these models for social scientific inquiry (which is what they were intended for) does not necessarily touch upon the humanistic and scientific dimensions of cultural and environmental topics.\footnote{\cite{Littlejohn2011a} outline three modes of scholarly inquiry: scientific, social scientific, and humanistic.} Multilevel approaches can play an important role. That said, further conceptual work is necessary before they are employed at the crossroads of intercultural and environmental research. 

\section{Research Problem}

To understand intercultural aspects of environmental issues there is a need to add depth and scope (both conceptually and methodologically) to the existing research. Holistic approaches, such as the socio-ecological model, appear well suited to provide the necessary scope, since they integrate multiple levels of social organization. However, these models are often intended for social-scientific inquiry and do not account for the humanistic and natural scientific aspects of culture and ecology, respectively. 

In response, this dissertation aims to develop a methodological and conceptual framework for the analysis of environmental communication. It seeks to propose philosophical concepts that can serve as a basis for grasping the complexities of human communication in the context of ecological themes and issues. Moreover, these concepts need to be based on real-world data. 

\textbf{Research Question}

Based on the analysis of corpus data, what conceptual principles can guide the study of communication about natural resources and the environment? How do these principles apply to intercultural communication scholarship?

\section{Aims \& Structure of the Dissertation}

In addition to the primary aim of the dissertation (that is, to address the research problem and question stated above), there are several complementary aims. In this dissertation, corpus linguistic methods are employed to collect and analyse data. One aim is to outline how corpus-linguistic methodologies can apply to both ecological and intercultural communication research. As elaborated in the methodology (Chapter 2), corpus linguistics is a powerful approach to studying communication, but is relatively uncommon in environmental and, to a lesser extent, intercultural communication studies. 

While the research problem is conceptual, the aim is to root concepts in real world data. So, rather than begin with a theoretical framework and proceed to methodology and results, this dissertation begins with data analysis and concludes with a framework. Although the focus is environmental, the dissertation addresses intercultural communication research more broadly, albeit in ways that are somewhat unconventional in the discipline.     

This dissertation is structured as follows: Chapter 2 introduces multilevel analysis in the context of intercultural communication research. This chapter also introduces corpus linguistics as a methodology to address the research problem. The multiple levels of analysis are then applied to three separate linguistic corpora that contain data on different environmental themes as well as different types of communication (i.e. textual, verbal, nonverbal). Chapters 3-5 are analyses of these three corpora, using the multilevel framework. Chapter, 6 integrates the analyses and develops conceptual principles to address the first part of the research question: ``what conceptual principles can guide the study of communication about natural resources and the environment?'' Finally, Chapter 7 addresses the second part of the research question (''How do these principles apply to intercultural communication scholarship?'') by discussing the findings in terms of current and future directions in intercultural communication research. 

Following the main chapters, there are three Appendices (A, B, and C) corresponding to chapters 3, 4, and 5, respectively. The appendices have links to the raw corpus data and well as programming code used for data processing and analysis.